%قبل از شروع مسئله ترجیحاً چیزی ننویسید. اما اگر خواستید بنویسید، در قالب لیستی (فقرات و شمارش) بنویسید.
%برای اضافه کردن نمره، بعدی از شروع‌مسئله، [x] را اضافه کنید و به جای x  نمره مورد نظر را قرار دهید. برای زیرمسئله‌ها هم به همین صورت است.
%\شروع‌مسئله[10]
\شروع‌مسئله

% صورت اصلی مسئله را در این‌جا بنویسید.

شیوه خاصی از تابع درهم‌سازی به روش تقسیم را در نظر بگیرید که در آن
$h(k) = k \operatorname{mod} m$
که در آن $m=2^p -1$ و $k$ هم یک رشته از کاراکترهاست که بیانگر عددی در مبنای $2^p$ است. نشان دهید که اگر دو رشته $x$ و $y$ داشته باشیم که با جابه‌جا کردن کاراکترهای یکی از آن‌ها بتوانیم به دیگری برسیم (یکی از آن‌ها جایگشتی از دیگری باشد)، هر دوی این رشته‌ها با تابع درهم‌سازی ذکر شده به یک خانه نگاشت می‌شوند. مثالی از یک کاربرد واقعی بیاورید که وجود چنین خاصیتی نامطلوب باشد.
\solution{
	نشان می‌دهیم که هش هر رشته در روش گفته شده برابر جمع ارقامش به پیمانه $2^p -1$ خواهد بود. برای این کار از استقرا استفاده می‌کنیم. برای پایه استقرا یک رشته تک کاراکتری را در نظر بگیرید. در این صورت مقدار آن به پیمانه $m=2^p -1$ محاسبه می‌شود. پس برای پایه استقرا حکم برقرار است. حال برای گام استقرا فرض کنید رشته $w=w_1 w_2$ باشد که از $n$ کاراکتر تشکیل شده است و $w_1$ شامل $n-1$ کاراکتر پرارزش تر باشد و $w_2$ هم کاراکتر آخر آن باشد. در این صورت
	$$h(w) = h(w_1) 2^p + h(w_2) \operatorname{mod} 2^p -1$$
	$$h(w) = h(w_1) + h(w_2)$$
	
	این یعنی  که برای یک رشته با طول دلخواه، حاصل تابع درهم ساز معادل حاصل جمع هش تمامی آن به جز رقم آخر و هش رقم آخر به صورت جداگانه است. این موضوع نشان می‌دهد که هش هر رقم به طور جداگانه تاثیر دارد و در نتیجه اگر یک عدد $k$ رقمی به فرم
	$x = x_1 x_2 \cdots x_k$
	داشته باشیم
	$$h(x) = h(x_k) + h(x_1 \cdots x_{k-1}) = h(x_k) + h(x_{k-1}) + h(x_1 \cdots x_{k-2}) = \cdots = h(x_1) + \cdots h(x_k)$$

خواهد بود.

برای مثال حالتی که چنین سیستمی مشکل ساز می‌شود، فرض کنید که یک دیکشنری انگلیسی طراحی می‌کنید که لغات براساس این تابع در آن ذخیره می‌شوند و کاراکترهای مورد استفاده ما هم کدهای ASCII باشند و در نتیجه مبنا $2^8 = 256$ و $m=255$ باشد. در این صورت مثلا کلماتی نظیر \lr{STOP,TOPS,SPOT,POTS} که همگی معنی دار هم هستند به یک slot نگاشت خواهند شد.

}



%%%%%%%%%%%%%%%%%%%%%%%%%%%%%%%%%%%%%%%%%%%%%%%%%%%%%
% برخی دستورات مفید برای تمرین‌ها. %
%%%%%%%%%%%%%%%%%%%%%%%%%%%%%%%%%%%%%%%%%%%%%%%%%%%%%

% فرمول‌های نمونه:
%$f(x) = \Theta(x)$
%$T(x, y) \leq \log(x) + 2^y + \binom{2n}{n}$
% {\tt code\_function}


% لیست موارد بدون شماره‌گذاری:
%\شروع{فقرات}
%    \فقره مورد اول 
%        \شروع{فقرات}
%            \فقره مورد اول داخلی 
%                \شروع{فقرات}
%                    \فقره اولین مورد خیلی داخلی 
%                    \فقره دومین مورد خیلی داخلی 
%                \پایان{فقرات}
%            \فقره مورد دوم داخلی 
%        \پایان{فقرات}
%    \فقره مورد دوم 
%\پایان{فقرات}

% لیست موارد شماره‌گذاری شده:
%\شروع{شمارش}
%    \فقره یک مورد 
%    \فقره یک مورد دیگر 
%    \فقره باز هم یک مورد دیگر
%\پایان{شمارش}

% برای نوشتن یک فرمول چند خطی که یک شماره بخوره. برای حالت‌های دیگر به سایت زیر مراجعه کنید:
% https://www.overleaf.com/learn/latex/Aligning_equations_with_amsmath
%\begin{align} 
%    \begin{split}
%        (x+y)^3 &= (x+y)^2(x+y) \\
%                &= (x^2+2xy+y^2)(x+y) \\
%                &= (x^3+2x^2y+xy^2) + (x^2y+2xy^2+y^3) \\
%                &= x^3+3x^2y+3xy^2+y^3
%    \end{split}                                 
%\end{align}

% نوشتن ماتریس و آرایه دوبعدی
%\begin{align}
%    A = 
%    \begin{bmatrix}
%        A_{11} & A_{21} \\
%        A_{21} & A_{22}
%    \end{bmatrix}
%\end{align}

% اضافه کردن تصویر با شرح. برای اضافه کردن تصویر بدون شرح، فقط دستور درج‌تصویر را استفاده کنید.
%\شروع{شکل}[h!]
%\تنظیم‌ازوسط
%  \درج‌تصویر[پهنا=0.5\پهنای‌متن]{img.jpg}
%\شرح{شرح تصویر}
%\برچسب{تصویر:۱}
%\پایان{شکل}
%\پایان{زیرمسائل}

%کد زیر را ببینید.
%%\lstinputlisting{source_filename.py}
%\begin{latin}
%\begin{lstlisting}
%# Import the modules
%import sys
%import random
%
%ans = True
%
%while ans:
%question = raw_input("Ask the magic 8 ball a question: (press enter to quit) ")
%
%answers = random.randint(1,8)
%
%if question == "":
%sys.exit()
%
%elif answers == 1:
%print "It is certain"
%
%elif answers == 2:
%print "Outlook good"
%
%elif answers == 3:
%print "You may rely on it"
%
%elif answers == 4:
%print "Ask again later"
%
%elif answers == 5:
%print "Concentrate and ask again"
%
%elif answers == 6:
%print "Reply hazy, try again"
%
%elif answers == 7:
%print "My reply is no"
%
%elif answers == 8:
%print "My sources say no"
%\end{lstlisting}
%\end{latin}
